\section*{Supplementary material}
\label{sec:appendix}

% ----------------------------- SUBSECTION  -----------------------------------
\subsection*{Proof of Theorem~\ref{thm:unidentifiability}}

\begin{theorem*}
  Given a joint distribution over the protected attribute~$A$, the true
  label~$Y$, and some features~$X_1, \dots, X_n$, in which we have already
  specified the resolving variables, no observational criterion can generally
  determine whether the Bayes optimal unconstrained predictor or the Bayes
  optimal equal odds predictor exhibit unresolved discrimination.
\end{theorem*}

\begin{proof}
  Let us consider the two graphs in \figref{fig:unidentifiability}. First, we
  show that these graphs can generate the same joint distribution~$\bP(A, Y,
  X_1, X_2, \hR^*)$ for the Bayes optimal unconstrained predictor~$\hR^*$.

  We choose the following structural equations for the graph on the
  left\footnote{$\sigma(x) = 1 / (1 + e^{-x})$}
  \begin{itemize}
    \item $A = Ber(\nicefrac{1}{2})$
    \item $X_1$ is a mixture of Gaussians $\cN(A+1, 1)$ with weight~$\sigma(2
      A)$ and $\cN(A-1, 1)$ with weight~$\sigma(-2 A)$
    \item $Y = Ber( \sigma(2 X_1) )$
    \item $X_2 = X_1 - A$
    \item $\hR^* =  X_1$
    \item ($\tilde{\hR} = X_2$)
  \end{itemize}
  where the Bernoulli distribution~$Ber(p)$ without a superscript has
  support~$\{-1, 1\}$.

  For the graph on the right, we define the structural equations
  \begin{itemize}
    \item $A = Ber(\nicefrac{1}{2})$
    \item $Y = Ber(\sigma(2 A))$
    \item $X_2 = \cN(Y,1)$
    \item $X_1 = A + X_2$
    \item $\hR^* = X_1$
    \item ($\tilde{\hR} = X_2$)
  \end{itemize}

  First we show that in both scenarios~$\hR^*$ is actually an optimal score.
  In the first scenario~$Y \indep A \given X_1$ and~$Y \indep X_2 \given X_1$
  thus the optimal predictor is only based on~$X_1$. We find
  \begin{equation}\label{eq:YgivenX1}
    \Pr(Y = y \given X_1 = x_1) = \sigma(2 x_1 y) \eqc
  \end{equation}
  which is monotonic in~$x_1$. Hence optimal classification is obtained by
  thresholding a score based only on~$\hR^* = X_1$.

  In the second scenario, because~$Y \indep X_1 \given \{A, X2\}$ the optimal
  predictor only depends on~$A, X_2$. We compute for the densities
  \begin{subequations}\label{eq:YgivenX2A}
    \begin{align}
      \bP(Y \given X_2, A) & = \frac{\bP(Y, X_2, A)}{\bP(X_2, A)} \\
      &= \frac{\bP(X_2, A \given Y) \bP(Y)}{\bP(X_2, A)} \\
      &= \frac{\bP(X_2 \given Y) \bP(A \given Y)
        \bP(Y)}{\bP(X_2, A)} \\
      &= \frac{\bP(X_2 \given Y) \frac{\bP(Y \given A)
        \bP(A)}{\bP(Y)} \bP(Y)}{\bP(X_2, A)} \\
      &= \frac{\bP(X_2 \given Y) \bP(Y \given A) \bP(A)}
        {\bP(X_2, A)} \eqc
    \end{align}
  \end{subequations}
  where for the third equal sign we use~$A \indep X_2 \given Y$. In the
  numerator we have
  \begin{equation}\label{eq:density}
    \bP(X_2 \given Y=y)(x_2) \bP(Y \given A=a)(y) \bP(A)(a)
    = f_{\cN(y,1)}(x_2) f_{Ber(\sigma(2 a))}(y) f_{Ber(\nicefrac{1}{2})}(a)\eqc
  \end{equation}
  where~$f_{D}$ is the probability density function of the distribution~$D$.
  The denominator can be computed by summing up~\eqref{eq:density} for~$y\in
  \{-1,1\}$. Overall this results in
  \begin{equation*}
    \Pr(Y = y \given X_2 = x_2, A=a) = \sigma(2 y (a + x_2)) \eqp
  \end{equation*}

  Since by construction~$X_1 = A + X_2$, the optimal predictor is again~$\hR^* =
  X_1$. If the joint distribution~$\bP(A, Y, \hR^*)$ is identical in the two
  scenarios, so are the joint distributions~$\bP(A, Y, X_1, X_2, \hR^*)$,
  because of~$X_1 = \hR^*$ and~$X_2 = X_1 - A$.

  To show that the joint distributions~$\bP(A, Y, \hR^*) = \bP(Y \given A,
  \hR^*) \bP(\hR^* \given A) \bP(A)$ are the same, we compare the conditional
  distributions in the factorization.

  Let us start with~$\bP(Y \given A, \hR^*)$.  Since~$\hR^* = X_1$ and in the
  first graph~$Y \indep A \given X_1$, we already found the distribution
  in~\eqref{eq:YgivenX1}. In the right graph,~$\bP(Y \given \hR^*, A) = \bP(Y
  \given X_2 + A, A) = \bP(Y \given X_2, A)$ which we have found
  in~\eqref{eq:YgivenX2A} and coincides with the conditional in the left graph
  because of~$X_1 = A + X_2$.

  Now consider~$\hR^* \given A$. In the left graph we have~$\bP(\hR^* \given A)
  = \bP(X_1 \given A)$ and the distribution~$\bP(X_1 \given A)$ is just the
  mixture of Gaussians defined in the structural equation model. In the right
  graph~$\hR^* = A + X_2 = Y + \cN(A,1)$ and thus~$\bP(\hR^* \given A) = \cN(A
  \pm 1)$ for~$Y = \pm 1$. Because of the definition of~$Y$ in the structural
  equations of the right graph, following a Bernoulli distribution with
  probability~$\sigma(2 A)$, this is the same mixture of Gaussians as the one we
  found for the left graph.

  Clearly the distribution of~$A$ is identical in both cases.

  Consequently the joint distributions agree.

  When~$X_1$ is an resolving variable, the optimal predictor in the left
  graph does not exhibit unresolved discrimination, whereas the graph on the
  right does.

  The proof for the equal odds predictor~$\tilde{\hR}$ is immediate once we
  show~$\tilde{\hR} = X_2$. This can be seen from the graph on the right,
  because here~$X_2 \indep A \given Y$ and both using~$A$ or~$X_1$ would
  violate the equal odds condition. Because the joint distribution in the left
  graph is the same,~$\tilde{\hR} = X_2$ is also the optimal equal odds
  score.\qedhere
\end{proof}

% ----------------------------- SUBSECTION  ------------------------------------
\subsection*{Proof of Proposition~\ref{pro:unawareness}}

\begin{proposition*}
  If there is no directed path from a proxy to a feature, unawareness avoids
  proxy discrimination.
\end{proposition*}

\begin{proof}
  An unaware predictor~${\hR}$ is given by~${\hR} = r (X)$ for some function~$r$
  and features~$X$. If there is no directed path from proxies~$P$ to~$X$,
  \ie{}~$P \notin ta^{\cG}(X)$, then~${\hR} = r(X) = r(ta^{\cG}(X)) =
  r(ta^{\cG}_P(X))$.  Thus~$\bP({\hR} \given do(P=p)) = \bP({\hR})$ for all~$p$,
  which avoids proxy discrimination.\qedhere
\end{proof}

% ----------------------------- SUBSECTION  ------------------------------------
\subsection*{Proof of Theorem~\ref{thm:rvtodistr}}

\begin{theorem*}
  Let the influence of~$P$ on~$X$ be additive and linear, \ie{}
  \begin{equation*}
    X = f_X (P, ta^{\cG}_P(X)) = g_X(ta^{\cG}_P(X)) + \mu_X P
  \end{equation*}
  for some function~$g_X$ and~$\mu_X \in \bR$. Then any predictor of the form
  \begin{equation*}
    \hR = r(X - \E [X \given do(P)])
  \end{equation*}
  for some function~$r$ exhibits no proxy discrimination.
\end{theorem*}

\begin{proof}
  It suffices to show that the argument of~$r$ is constant \wrt{} to~$P$,
  because then~$\hR$ and thus~$\bP(\hR)$ are invariant under changes of~$P$. We
  compute
  \begin{align*}
    \E[X \given do(P)] &= \E[ g_X(ta^{\cG}_P (X)) + \mu_X P \given do(P)] \\
    &= \usub{=0}{\E[ g_X(ta^{\cG}_P (X)) \given do(P)]} + \E[\mu_X P \given do(P)]\\
    &= \mu_X P \eqp
  \end{align*}
  Hence,
  \begin{equation*}
    X - \E[X \given do(P)] = g_X(ta^{\cG}_P (X))
  \end{equation*}
  is clearly constant \wrt{} to~$P$.\qedhere
\end{proof}

% ----------------------------- SUBSECTION  ------------------------------------
\subsection*{Proof of Corollary~\ref{cor:rvtodistr}}

\begin{corollary*}
  Under the assumptions of Theorem~\ref{thm:rvtodistr}, if all directed paths
  from any ancestor of~$P$ to~$X$ in the graph~$\cG$ are blocked by~$P$, then
  any predictor based on the \emph{adjusted features}~$\tilde{X} := X - \E[X
  \given P]$ exhibits no proxy discrimination and can be learned from the
  observational distribution~$\bP(P, X, Y)$ when target labels~$Y$ are
  available.
\end{corollary*}

\begin{proof}
  Let~$Z$ denote the set of ancestors of~$P$. Under the given assumptions~$Z
  \cap ta^{\cG} (X) = \emptyset$, because in~$\cG$ all arrows into~$P$ are
  removed, which breaks all directed paths from any variable in~$Z$ to~$X$ by
  assumption. Hence the distribution of~$X$ under an intervention on~$P$
  in~$\tcG$, where the influence of potential ancestors of~$P$ on~$X$ that
  does not go through~$P$ would not be affected, is the same as simply
  conditioning on~$P$. Therefore~$\E[X \given do(P)] = \E[X \given P]$, which
  can be computed from the joint observational distribution, since we
  observe~$X$ and~$P$ as generated by~$\tcG$.\qedhere
\end{proof}

% ----------------------------- SUBSECTION  ------------------------------------
\subsection*{Proof of Proposition~\ref{pro:inexpectation}}

\begin{proposition*}
  Any predictor of the form~$\hR = \lambda (X - \E[X \given do(P)]) + c$ for
  linear~$\lambda, c \in \bR$ exhibits no proxy discrimination in expectation.
\end{proposition*}

\begin{proof}
  We directly test the definition of proxy discrimination in expectation using
  the linearity of the expectation
  \begin{align*}
    \E[\hR \given do(P=p)] &= \E[\lambda (X - \E[X \given do(P)])
      + c \given do(P=p)] \\
    &= \lambda (\E[ X \given do(P=p)] - \E[X \given do(P=p)]) + c \\
    &= c \eqp
  \end{align*}
  This holds for any~$p$, hence proxy discrimination in expectation is
  achieved.\qedhere
\end{proof}

% ----------------------------- SUBSECTION  -----------------------------------
\subsection*{Additional statements}

Here we provide an additional statement that is a first step towards the
``opposite direction'' of Theorem~\ref{thm:rvtodistr}, \ie{} whether we can
infer information about the structural equations, when we are given a predictor
of a special form that does not exhibit proxy discrimination.

\begin{theorem*}
  Let the influence of~$P$ on~$X$ be additive and linear and let the influence
  of~$P$ on the argument of~$\hR$ be additive linear, \ie{}
  \begin{align*}
    f_X (ta^{\cG}(X)) &= g_X(ta^{\cG}_P(X)) + \mu_X P \\
    f_{\hR} (P, ta^{\cG}_P(X)) &= h(g_{\hR}(ta^{\cG}_P(X)) + \mu_{\hR} P)
  \end{align*}
  for some functions~$g_X, g_{\hR}$, real numbers~$\mu_X, \mu_{\hR}$ and a
  smooth, strictly monotonic function~$h$. Then any predictor that avoids proxy
  discrimination is of the form
  \begin{equation*}
    {\hR} = r(X - \E[X \given do(P)])
  \end{equation*}
  for some function~$r$.
\end{theorem*}

\begin{proof}
  From the linearity assumptions we conclude that
  \begin{equation*}
    \hat{f}_{\hR}(P, X) = h(g_X(ta^{\cG}_P(X)) + \mu_X P + \hat{\mu}_{\hR} P) \eqc
  \end{equation*}
  with~$\hat{\mu}_{\hR} = \mu_{\hR} - \mu_P$ and thus~$g_X = g_{\hR}$. That
  means that both the dependence of~$X$ on~$P$ along the path~$P \to \dots \to
  X$ as well as the direct dependence of~${\hR}$ on~$P$ along~$P \to {\hR}$
  are additive and linear.

  To avoid proxy discrimination, we need
  \begin{subequations}\label{eq:distreq}
  \begin{align}
    \bP ({\hR} \given do(P=p))
    &= \bP(h(g_{\hR}(ta^{\cG}_P(X)) + \mu_{\hR} p)) \\
    &\overset{!}{=} \bP(h(g_{\hR}(ta^{\cG}_P(X)) + \mu_{\hR} p'))
    = \bP ({\hR} \given do(P=p')) \eqp
  \end{align}
  \end{subequations}

  Because~$h$ is smooth an strictly monotonic, we can conclude that already the
  distributions of the argument of~$h$ must be equal, otherwise the
  transformation of random variables could not result in equal distributions,
  \ie{}
  \begin{equation*}
    \bP(g_{\hR}(ta^{\cG}_P(X)) + \mu_{\hR} p)
    \overset{!}{=} \bP(g_{\hR}(ta^{\cG}_P(X)) + \mu_{\hR} p') \eqp
  \end{equation*}
  Since, up to an additive constant, we are comparing the distributions of the
  \emph{same} random variable~$g_{\hR}(ta^{\cG}_P(X))$ and not merely
  identically distributed ones, the following condition is not only
  sufficient, but also necessary for~\eqref{eq:distreq}
  \begin{equation*}
    g_{\hR}(ta^{\cG}_P(X)) + \mu_{\hR} p
    \overset{!}{=} g_{\hR}(ta^{\cG}_P(X)) + \mu_{\hR} p' \eqp
  \end{equation*}
  This holds true for all~$p, p'$ only if~$\mu_{\hR} = 0$, which is equivalent
  to~$\hat{\mu}_{\hR} = - \mu_P$.

  Because as in the proof of~\ref{thm:rvtodistr}
  \begin{equation*}
    \E[X \given do(P)] = \mu_X P,
  \end{equation*}
  under the given assumptions any predictor that avoids proxy discrimination is
  simply
  \begin{equation*}
    {\hR} = X + \mu_{\hR} P = X - \E[X \given do(P)] \eqp
  \end{equation*}\qedhere
\end{proof}
